\documentclass[11pt]{article}
\usepackage[margin=1in]{geometry}
\usepackage{amsmath, amssymb}
\usepackage{graphicx}
\usepackage{natbib}
\usepackage{booktabs}
\usepackage{hyperref}
\usepackage{xcolor}
\usepackage{authblk}
\usepackage[section]{placeins}  % keep floats within their section (auto \FloatBarrier before each \section)

\newcommand{\Msun}{M_\odot}
\newcommand{\code}[1]{\texttt{#1}}
\newcommand{\todo}[1]{\textcolor{red}{[TODO: #1]}}

\title{Differentiable Forward Models and Hierarchical Inference for\\
Weak-Lensing Galaxy Cluster Mass Calibration:\\
A Comparison of HMC and Simulation-Based Inference}

\usepackage{orcidlink}

\author[1,2]{Akum Gill\,\orcidlink{0000-0002-4383-4860}}
\author[3,4]{Camille Avestruz\,\orcidlink{0000-0001-8868-0810}}
\author[3,4]{Ming-Feng Ho\,\orcidlink{0000-0002-4457-890X}}
\affil[1]{\small Department of Physics, Harvard University}
\affil[2]{\small Center for Astrophysics \textbar\ Harvard \& Smithsonian}
\affil[3]{\small Department of Physics, University of Michigan}
\affil[4]{\small Leinweber Institute for Theoretical Physics, University of Michigan}
\date{Draft --- \today}

\begin{document}
\maketitle

\begin{abstract}
In a companion paper we showed that simulation-based inference (SBI) recovers the intrinsic
population distribution of galaxy-cluster mass and concentration from stacked weak-lensing
profiles, whereas a standard Markov Chain Monte Carlo (MCMC) joint-likelihood fit recovers only a
central value and therefore appears ``overconfident'' relative to the population spread. That
comparison left open a natural question: is the limitation fundamental to likelihood-based
inference, or merely to the \emph{single-parameter} model the joint likelihood assumes? Here we
answer it directly. We build a fully differentiable NFW forward model in \code{JAX} (using
\code{halox}) and implement a \emph{hierarchical} Bayesian model that infers per-cluster
$(M,c)$ jointly with population-level hyperparameters, sampled with the No-U-Turn Sampler (NUTS).
Across all of our numerical experiments the hierarchical model recovers the true population
mass spread $\sigma_{\log_{10}M}$ to within a few percent --- matching SBI and resolving the
apparent MCMC overconfidence --- while remaining computationally tractable
($\sim$3--5\,min per experiment on a single CPU). We further identify a noise-induced
concentration bias inherent to analytic-likelihood inference that SBI avoids by construction, and
we discuss the practical trade-offs (amortization, differentiability, convergence diagnostics)
between SBI and hierarchical HMC for upcoming surveys such as LSST.
\end{abstract}

\section{Introduction}
\label{sec:intro}

Weak-lensing mass calibration underpins cluster cosmology: the cluster mass function is a sensitive
probe of structure growth, but only if cluster masses (and their scaling with observables such as
optical richness) are accurately and unbiasedly recovered \citep{Pratt2019}. Surveys stack the
weak-lensing signal of many clusters to beat down shape noise, and then infer mass (and, in our
framework, concentration) from the stacked profile or the ensemble of individual profiles.

The companion paper \citep[][hereafter Paper~I]{paper1} compared two inference engines --- MCMC and
SBI --- and two ways of aggregating clusters: \emph{join-then-fit} (JTF; fit the stacked median
profile) and \emph{fit-then-join} (FTJ; a joint analysis of all profiles). A central result of
Paper~I is that the MCMC FTJ posterior is far narrower than the true population distribution of
$(M,c)$, because the joint likelihood treats all $N_c$ clusters as repeated measurements of a
\emph{single} underlying $(M,c)$ pair. Its posterior therefore estimates a central
(approximately mean) value, not the population spread --- it answers a different question than the
SBI FTJ inferrer, which is trained to output population percentiles.

Paper~I noted that recovering the population spread with a likelihood-based method ``would require
a hierarchical model that explicitly infers per-cluster $(M,c)$ together with population-level
hyperparameters.'' This paper carries out exactly that program. Our contributions are:
\begin{enumerate}
  \item A differentiable NFW surface-density forward model in \code{JAX}/\code{halox}
        (Section~\ref{sec:forward}), validated against the \code{Colossus}-based pipeline used to
        generate the data.
  \item A hierarchical Bayesian model (Section~\ref{sec:hbm}) sampled with NUTS that recovers the
        population $(M,c)$ distribution, closing the gap with SBI (Section~\ref{sec:results}).
  \item A three-way comparison (SBI, MCMC joint-likelihood, hierarchical HMC) of population-spread
        recovery and computational cost across all of Paper~I's experiments
        (Section~\ref{sec:results}, Table~\ref{tab:three_method}).
  \item Identification of a noise-induced concentration bias intrinsic to analytic-likelihood
        inference, which SBI avoids (Section~\ref{sec:noisebias}).
\end{enumerate}

\section{Differentiable forward model}
\label{sec:forward}

Hamiltonian Monte Carlo and its adaptive variant NUTS require gradients of the log-posterior with
respect to all parameters. The \code{Colossus} \citep{Colossus18} NFW routines used to generate our
mock data are not differentiable, so we re-express the forward model in \code{JAX}. We use
\code{halox} \citep{halox} for the (differentiable) NFW scale radius $R_s$ and characteristic
density $\rho_0$ given $(M_{\rm vir}, c, z)$ in the Planck18 cosmology, and evaluate the projected
surface density with the analytic NFW expression of \citet{Bartelmann1996} / \citet{WrightBrainerd2000},
\begin{equation}
  \Sigma(R) = 2\,\rho_0 R_s\, f(R/R_s),
\end{equation}
with the standard piecewise $f(x)$. We implement $f$ with a ``safe-$x$'' guard so that the
untaken branch of the piecewise expression never produces a NaN gradient. The result matches
\code{halox}'s own surface-density routine to machine precision and the \code{Colossus} data
generator to better than $0.1\%$ across all radial bins, while remaining differentiable
(verified by finite, well-behaved gradients $\partial \log_{10}\Sigma/\partial \log_{10}M$ and
$\partial \log_{10}\Sigma/\partial c$).

\section{Hierarchical Bayesian model}
\label{sec:hbm}

We model a population of $N_c$ clusters in a fixed richness bin. The generative model is
\begin{align}
  \log_{10}M_j &\sim \mathcal{N}(\mu_M,\ \sigma_M), & j &= 1,\dots,N_c,\\
  c_j &\sim \mathcal{N}\!\big(c_0 + \beta\,(\log_{10}M_j - \mu_M),\ \sigma_c\big),\\
  d_{j} &\sim \mathcal{N}\!\big(\log_{10}\Sigma(\log_{10}M_j, c_j),\ \sigma_{\rm obs}\big),
\end{align}
where $(\mu_M,\sigma_M)$ are the population mean and intrinsic spread in log mass; the
mass--concentration relation has intercept $c_0$, slope $\beta$, and intrinsic scatter $\sigma_c$;
and $d_j$ is the observed log-surface-density profile of cluster $j$. We additionally infer a
fractional log-scatter term $\sigma_{\rm extra}$ added in quadrature to $\sigma_{\rm obs}$,
the hierarchical analogue of the error-inflation parameter used in the analytic-likelihood
analyses of Paper~I. The per-cluster latents $(\log_{10}M_j, c_j)$ are sampled in a non-centred
parameterization to avoid the funnel pathology characteristic of hierarchical models
\citep{Betancourt2015}.

The key structural difference from the FTJ joint likelihood of Paper~I is that here the
per-cluster parameters are drawn from a population with \emph{free} width $\sigma_M, \sigma_c$;
the joint likelihood is the special case $\sigma_M \equiv 0$ (one shared $(M,c)$). The posterior on
$(\sigma_M, \sigma_c)$ is therefore a direct estimate of the population spread --- precisely the
quantity that SBI FTJ is trained to output and that the joint likelihood cannot represent.
This structural distinction is made explicit in Figures~\ref{fig:pgm_simple} and~\ref{fig:pgm_hier}.

\begin{figure}[t]
  \centering
  \includegraphics[width=\columnwidth]{figures/pgm_jtf_ftj.png}
  \caption{Probabilistic graphical models for the two single-$(M,c)$ inference structures of
  Paper~I, generated directly from the model code with \code{numpyro.render\_model}. Shaded nodes
  are observed; the box (``plate'') denotes replication over the $N_c$ clusters. \textit{Left,
  join-then-fit:} a single $(\log_{10}M, c)$ generates the one stacked (median) profile.
  \textit{Right, fit-then-join:} a single $(\log_{10}M, c)$ sits \emph{outside} the plate and
  generates all $N_c$ observed profiles---one shared $(M,c)$ is assumed for every cluster, so the
  posterior estimates a central value, not the population spread. ($\ln f$ is the shared
  error-inflation nuisance parameter.)}
  \label{fig:pgm_simple}
\end{figure}

\begin{figure}[t]
  \centering
  \includegraphics[width=0.95\columnwidth]{figures/pgm_hier_solo.png}
  \caption{Probabilistic graphical model for the hierarchical approach (this work), generated with
  \code{numpyro.render\_model}. In contrast to the single-$(M,c)$ models of
  Figure~\ref{fig:pgm_simple}, population hyperparameters
  ($\mu_M, \sigma_M, c_0, \beta, \sigma_c$) generate \emph{per-cluster} latent parameters
  ($\log_{10}M_j, c_j$) \emph{inside} the plate, which in turn generate the observed profiles
  $\Sigma_j$. This additional layer of per-cluster latents, drawn from a population of free width,
  is exactly what allows the model to infer the intrinsic population spread.}
  \label{fig:pgm_hier}
\end{figure}

We sample with NUTS \citep{Hoffman2014} as implemented in \code{NumPyro} \citep{numpyro}, using
four chains, $1000$ warmup and $1500$ sampling iterations each, and a target acceptance
probability of $0.9$. Convergence is assessed with the split-$\hat r$ statistic and the effective
sample size; all reported fits satisfy $\hat r < 1.01$.

\section{Results}
\label{sec:results}

Throughout this section, ``SBI'' (or ``SBI FTJ'') refers to the percentile-summary fit-then-join
inferrer of Paper~I---trained to output percentiles of the population $(M,c)$ distribution, which
are then fit by a Gaussian---unless explicitly qualified as ``amortized hierarchical SBI,'' the
distinct approach of Section~\ref{sec:amortized_hbi}. All three-method comparisons (Table~\ref{tab:three_method},
Figures~\ref{fig:aggregate}, \ref{fig:calibration}, \ref{fig:ppc}) use the Paper~I SBI FTJ.

\begin{table*}
\centering
\caption{Recovered population \emph{spread} ($\sigma_{\log_{10}M}$, $\sigma_c$) and \emph{center} ($\mu_{\log_{10}M}$, $\mu_c$) across experiments, for the three inference approaches, compared to the true population. For the spread, \textbf{bold} marks recovery within 20\% of truth and $^{\dagger}$ marks a value off by more than a factor of two (strongly over- or under-confident); for the center, \textbf{bold} marks recovery within $0.03$ (mass) or $0.2$ (concentration) of truth. MCMC joint-likelihood (FTJ) reports its posterior \emph{width} as the implied spread, which collapses far below the truth because it estimates a central value, not a population; SBI FTJ and hierarchical HMC report the inferred population. All HMC fits converged ($\hat r<1.01$).}
\label{tab:three_method}
\resizebox{\textwidth}{!}{%
\begin{tabular}{l rrrr rrrr}
\hline
\multicolumn{9}{l}{\textit{Population spread}} \\
 & \multicolumn{4}{c}{$\sigma_{\log_{10}M}$} & \multicolumn{4}{c}{$\sigma_c$} \\
Experiment & TRUE & MCMC & SBI & HMC & TRUE & MCMC & SBI & HMC \\
\hline
baseline & 0.118 & 0.006$^{\dagger}$ & \textbf{0.115} & \textbf{0.117} & 0.178 & 0.048$^{\dagger}$ & \textbf{0.189} & \textbf{0.166} \\
baseline high mc scatter & 0.118 & 0.006$^{\dagger}$ & \textbf{0.111} & \textbf{0.116} & 0.719 & 0.048$^{\dagger}$ & 0.187$^{\dagger}$ & \textbf{0.794} \\
baseline high noise & 0.118 & 0.013$^{\dagger}$ & \textbf{0.119} & \textbf{0.127} & 0.178 & 0.080$^{\dagger}$ & \textbf{0.191} & 0.276 \\
baseline high richness contam & 0.162 & 0.006$^{\dagger}$ & \textbf{0.140} & \textbf{0.154} & 0.212 & 0.048$^{\dagger}$ & \textbf{0.206} & \textbf{0.217} \\
baseline high rm scatter & 0.535 & 0.000$^{\dagger}$ & 0.208$^{\dagger}$ & \textbf{0.517} & 0.739 & 0.049$^{\dagger}$ & 0.301$^{\dagger}$ & 0.103$^{\dagger}$ \\
baseline low richness contam & 0.133 & 0.006$^{\dagger}$ & \textbf{0.133} & \textbf{0.136} & 0.186 & 0.047$^{\dagger}$ & \textbf{0.189} & \textbf{0.182} \\
baseline ludlow & 0.118 & 0.006$^{\dagger}$ & \textbf{0.118} & \textbf{0.117} & 0.219 & 0.049$^{\dagger}$ & 0.173 & 0.175 \\
baseline prada & 0.118 & 0.006$^{\dagger}$ & 0.142 & \textbf{0.118} & 0.137 & 0.052$^{\dagger}$ & \textbf{0.152} & 0.207 \\
\hline
\multicolumn{9}{l}{\textit{Population center}} \\
 & \multicolumn{4}{c}{$\mu_{\log_{10}M}$} & \multicolumn{4}{c}{$\mu_c$} \\
Experiment & TRUE & MCMC & SBI & HMC & TRUE & MCMC & SBI & HMC \\
\hline
baseline & 14.376 & \textbf{14.36} & \textbf{14.37} & \textbf{14.36} & 4.66 & \textbf{4.779} & \textbf{4.667} & \textbf{4.818} \\
baseline high mc scatter & 14.376 & \textbf{14.36} & \textbf{14.37} & \textbf{14.36} & 4.70 & \textbf{4.776} & \textbf{4.668} & \textbf{4.866} \\
baseline high noise & 14.376 & \textbf{14.36} & \textbf{14.35} & \textbf{14.35} & 4.66 & \textbf{4.814} & \textbf{4.676} & 5.056 \\
baseline high richness contam & 14.255 & \textbf{14.26} & \textbf{14.28} & \textbf{14.26} & 4.77 & \textbf{4.727} & \textbf{4.762} & \textbf{4.732} \\
baseline high rm scatter & 13.213 & 13.69 & 13.17 & \textbf{13.21} & 6.04 & 3.910 & \textbf{5.937} & \textbf{6.203} \\
baseline low richness contam & 14.353 & \textbf{14.35} & \textbf{14.38} & \textbf{14.35} & 4.68 & \textbf{4.755} & \textbf{4.656} & \textbf{4.787} \\
baseline ludlow & 14.376 & \textbf{14.37} & 14.41 & \textbf{14.36} & 5.30 & \textbf{5.271} & 4.658 & \textbf{5.473} \\
baseline prada & 14.376 & \textbf{14.39} & 14.49 & \textbf{14.36} & 6.15 & 5.862 & 4.639 & \textbf{6.336} \\
\hline
\end{tabular}
}
\end{table*}

Table~\ref{tab:three_method} summarizes the central result. In the
baseline (near-ideal) configuration, the true population mass spread is
$\sigma_{\log_{10}M} = 0.118$; the hierarchical HMC recovers $0.117$ and SBI FTJ recovers $0.116$,
while the MCMC joint-likelihood posterior width is $0.006$ --- roughly twenty times too small,
exactly the ``overconfidence'' identified in Paper~I. This pattern holds across experiments: the
hierarchical model recovers $\sigma_{\log_{10}M}$ to within a few percent of truth in every case,
including the high richness--mass scatter experiment ($0.52$ recovered vs.\ $0.54$ true), where
SBI underestimates the spread ($0.21$) because its training distribution does not span the
inflated scatter.

The hierarchical HMC posterior also reproduces the full two-dimensional $(M,c)$ population ellipse,
including the mass--concentration anti-correlation, as shown in the left column of
Figure~\ref{fig:aggregate}. We emphasize that this is achieved by a likelihood-based method:
the apparent overconfidence of MCMC in Paper~I was a property of the single-$(M,c)$ model, not of
likelihood-based inference per se.

\begin{figure*}[t]
  \centering
  \includegraphics[width=0.62\textwidth]{figures/aggregate_hbi_comparison.png}
  \caption{Aggregate population-recovery comparison across all experiments (one row each),
  analogous to Fig.~12 of Paper~I but adding the hierarchical HMC. Left: the $2\sigma$ $(M,c)$
  contour for each method against the true clusters (grey points, black-dashed true contour).
  Middle and right: the mass and concentration marginals. The MCMC joint-likelihood (blue)
  collapses to a narrow spike in every row (it estimates a central value); SBI FTJ (orange) and the
  hierarchical HMC (green) recover the true population width. The hierarchical model additionally
  captures the broad true distribution in the high $\lambda$--M scatter experiment that SBI FTJ
  underestimates.}
  \label{fig:aggregate}
\end{figure*}

Figure~\ref{fig:aggregate} extends this comparison to all experiments in the multi-panel format of
Paper~I, confirming that the hierarchical HMC tracks the true population marginals in both mass and
concentration across the suite.

\subsection{Calibration}
\label{sec:calibration}

\begin{figure}[t]
  \centering
  \includegraphics[width=0.82\columnwidth]{figures/calibration_three_method.png}
  \caption{Calibration of the three inference approaches for the baseline experiment. The
  horizontal axis is the nominal credible level $p$; the vertical axis is the empirical coverage,
  the fraction of the $N_c=376$ true cluster $(M,c)$ values that fall within the $p$ credible
  region of each inferred posterior (computed from the Mahalanobis distance to the posterior mean
  and covariance, so it accounts for the $M$--$c$ correlation). Perfect calibration lies on the
  diagonal; curves below it are overconfident. We summarize each with the maximum deviation from
  the diagonal, $\Delta_{\max}$. The MCMC joint-likelihood (blue, $\Delta_{\max}=0.93$) is
  severely overconfident---its narrow single-$(M,c)$ posterior contains almost none of the true
  population until $p\to1$---whereas SBI FTJ (orange, $\Delta_{\max}=0.02$) and the hierarchical
  HMC (green, $\Delta_{\max}=0.11$) are both well calibrated.}
  \label{fig:calibration}
\end{figure}

Unlike the Kolmogorov--Smirnov comparison of Paper~I, the coverage in
Figure~\ref{fig:calibration} is sample-size independent and accounts for the full $M$--$c$
covariance, making it our most robust calibration diagnostic. It confirms the headline result
directly: the single-$(M,c)$ MCMC posterior is dramatically overconfident
($\Delta_{\max}=0.93$), while both the hierarchical HMC and SBI FTJ recover near-nominal coverage
($\Delta_{\max}=0.11$ and $0.02$).

\begin{figure*}[t]
  \centering
  \includegraphics[width=0.95\textwidth]{figures/calibration_grid.png}
  \caption{Calibration across all experiments (fit-then-join), one panel each, in the format of
  Figure~\ref{fig:calibration}; the inset lists $\Delta_{\max}$ for MCMC (M), SBI (S) and
  hierarchical HMC (H). The MCMC joint-likelihood (blue) is severely overconfident in every
  in-distribution experiment. SBI FTJ (orange) and the hierarchical HMC (green) are both well calibrated
  in-distribution, with two informative exceptions: under strong M-c scatter SBI FTJ under-covers
  (its concentration spread is underestimated; cf.\ Table~\ref{tab:three_method}), and in the
  out-of-distribution cases (Prada, Ludlow) where the assumed scaling relation is wrong, all methods
  degrade. The hierarchical HMC is the most consistently calibrated across the in-distribution
  suite.}
  \label{fig:calibration_grid}
\end{figure*}

Figure~\ref{fig:calibration_grid} extends the comparison to all experiments. The pattern is
consistent: MCMC is overconfident throughout; SBI FTJ and the hierarchical HMC are well calibrated when
the model matches the data and degrade together under genuine model misspecification.

\subsection{Posterior predictive checks}
\label{sec:ppc}

\begin{figure*}[t]
  \centering
  \includegraphics[width=0.9\textwidth]{figures/aggregate_ppc_three_method.png}
  \caption{Three-method posterior predictive checks, one row per experiment. For each method
  (MCMC joint-likelihood, blue; SBI FTJ, orange; hierarchical HMC, green) we draw a population from
  its inferred $(\log_{10}M, c)$ distribution, forward-model each draw to a noiseless NFW profile,
  and compare to the observed profiles. \emph{Left:} grey lines are the observed profiles, the solid
  black line their median; each method shows its predicted-population $16$--$84$\% band and median.
  \emph{Right:} fractional residual of each method's predicted median relative to the observed
  median, $\Sigma_{\rm pred}/\Sigma_{\rm obs}-1$, with the method population bands. The predicted
  medians of all three methods reproduce the observed median in most experiments (the population
  \emph{mean} is well constrained by all); the methods differ in their predicted population
  \emph{spread} (band width) and in the out-of-distribution cases (Prada, Ludlow), where the
  residuals diverge at large radius.}
  \label{fig:ppc}
\end{figure*}

As an absolute goodness-of-fit check (rather than a comparison against the known truth, which would
be unavailable for real data), Figure~\ref{fig:ppc} shows posterior predictive checks for all three
methods: profiles forward-modelled from each inferred population, overlaid on the observed profiles.
The predicted population means reproduce the observed median across experiments, while the predicted
population spreads differ between methods consistent with Figure~\ref{fig:aggregate}.

\subsection{Jointly inferring the measurement noise}
\label{sec:infernoise}

The results above assume the per-bin measurement noise is known. A strength of the hierarchical
approach is that this assumption can be relaxed: we can promote the noise to free hyperparameters
and infer it jointly with the population. We replace the fixed per-bin uncertainties with an
independent free noise amplitude in each of the $30$ radial bins (a weakly-informative
$\mathrm{HalfNormal}(0.5)$ prior on each), and re-run the baseline fit ($\hat r<1.01$). Two
outcomes are notable (Figure~\ref{fig:infernoise}). First, the population mass spread is still
recovered accurately ($\sigma_M=0.124$ vs.\ true $0.118$), so freeing $30$ noise parameters does
\emph{not} destroy the population inference: the radial \emph{shape} of the signal---white per-bin
noise versus coherent NFW-shape population scatter---breaks the degeneracy. Second, the inferred
per-bin noise tracks the true input noise bin-by-bin (mean recovered amplitude $0.301$ vs.\ true
$0.307$\,dex; per-bin RMS residual $0.009$\,dex), demonstrating that the model recovers the
measurement uncertainty from the data alone. This capability---joint inference of population
properties and the noise model---is unavailable to the single-$(M,c)$ MCMC or the
percentile-summary SBI of Paper~I, and is a concrete advantage of the fully hierarchical
formulation.

\begin{figure}[t]
  \centering
  \includegraphics[width=0.9\columnwidth]{figures/inferred_noise.png}
  \caption{Jointly inferred per-bin profile noise (green; median and $16$--$84$\% band) compared to
  the true input noise used to generate the data (black dashed). With the noise promoted to $30$
  free hyperparameters, the model recovers both the per-bin measurement noise (RMS residual
  $0.009$\,dex) and the population mass spread ($\sigma_M=0.124$ vs.\ true $0.118$) from the data
  alone.}
  \label{fig:infernoise}
\end{figure}

The flat input noise above is a simplification; realistic stacked weak-lensing profiles have
spatially varying uncertainty, typically elevated at the smallest radii (deblending, miscentering,
and few background sources near the cluster centre) and at the largest radii (declining
signal-to-noise). To test this regime we regenerate the baseline population with a U-shaped per-bin
noise profile---$\sim$0.55\,dex at the innermost and outermost radii, falling to $\sim$0.20\,dex in
the middle---and again infer the $30$ per-bin noise amplitudes jointly with the population
(Figure~\ref{fig:ushaped}). The hierarchical model recovers the full U-shaped noise profile
bin-by-bin (RMS residual $0.011$\,dex) while still recovering the population mass spread
($\sigma_M=0.121$ vs.\ true $0.118$). This demonstrates that the joint population-plus-noise
inference is robust to realistic heteroscedastic noise, not merely the idealized homoscedastic
case.

\begin{figure*}[t]
  \centering
  \includegraphics[width=0.8\textwidth]{figures/ushaped_noise.png}
  \caption{Hierarchical HMC under realistic radially-varying (U-shaped) noise. \textit{Left:} the
  inferred per-bin noise (green, median and $16$--$84$\% band) recovers the true U-shaped input
  profile (black dashed)---high at the innermost and outermost radii, low in the middle---to an RMS
  residual of $0.011$\,dex. \textit{Right:} despite the complex noise, the recovered population mass
  marginal matches the truth ($\sigma_M=0.121$ vs.\ true $0.118$).}
  \label{fig:ushaped}
\end{figure*}

\subsection{A noise-induced concentration bias}
\label{sec:noisebias}

Both analytic-likelihood methods (MCMC joint and, more weakly, hierarchical HMC) center the
concentration high relative to the truth ($c \approx 4.8$ vs.\ $4.66$ in the baseline), whereas SBI
is unbiased. A diagnostic ladder isolates the cause. Fitting a single $(M,c)$ to the
\emph{noiseless} mean log-profile recovers the truth to $<0.001$; the hierarchical model on
noiseless data recovers $c$ to $+0.004$; only with realistic $0.3\,$dex profile noise does the
bias appear ($+0.15$). The bias is therefore noise-induced: per-bin noise combined with the
nonlinear NFW model and a Gaussian log-likelihood biases the recovered mass slightly low, and the
strong $M$--$c$ anti-correlation maps this into a concentration pulled high. SBI avoids the bias
because it learns the inverse mapping directly from \emph{noisy} training pairs, rather than
assuming the idealized analytic likelihood. This is a second, independent advantage of SBI
even for in-distribution inference, and a caution for analytic-likelihood cluster lensing analyses
operating at low signal-to-noise.

\subsection{Computational cost}
\label{sec:cost}

\begin{figure}[t]
  \centering
  \includegraphics[width=0.7\columnwidth]{figures/three_method_timing.png}
  \caption{Recurring inference cost grouped by task (join-then-fit, fit-then-join, and the full
  hierarchical population model) for the baseline experiment ($N_c=376$, single CPU), with bars
  ordered by increasing cost (SBI, HMC, \code{emcee}) within each group. For the single-$(M,c)$
  tasks, amortized SBI evaluation takes $\lesssim$1\,s, gradient-based HMC a few seconds, and the
  affine-invariant ensemble sampler (\code{emcee}) tens to thousands of seconds. For the
  $\sim$750-dimensional hierarchical model, amortized SBI runs in $\sim$10\,s and HMC in
  $\sim$210\,s; the hatched bar is an estimate for sampling the \emph{same} hierarchical model with
  \code{emcee}, which at this dimensionality requires $\gtrsim$1500 walkers (measured $\sim$1\,s per
  step, acceptance fraction $\sim$0.1) and is estimated to take of order a day or more---two to
  three orders of magnitude slower than gradient-based HMC. Across all tasks, gradient-based and
  amortized methods are dramatically cheaper than the gradient-free ensemble sampler, and the gap
  widens with dimensionality; SBI additionally carries a one-time training cost of $\sim$465\,s
  (excluded here).}
  \label{fig:timing}
\end{figure}

The differentiable forward model makes hierarchical HMC practical. Each experiment samples in
$\sim$3--5\,min on a single CPU core (four chains, $\hat r < 1.01$); running our full suite of
eight distinct observation sets costs $\sim$30\,min of sampling. Notably, this is roughly an order
of magnitude \emph{faster} than the affine-invariant ensemble MCMC of Paper~I once the latter is
run to the same autocorrelation-convergence standard ($\sim$2700\,s vs.\ $\sim$210\,s for the
baseline; Figure~\ref{fig:timing}), even though the hierarchical model infers $\sim$750 latent
parameters rather than three---gradient-based sampling scales far better in dimension than the
ensemble sampler, which mixes slowly on the strongly correlated joint likelihood. Indeed, applying
the affine-invariant ensemble sampler directly to the $\sim$750-dimensional hierarchical model is
estimated to take of order a day or more to converge (Figure~\ref{fig:timing}): it requires more
walkers than parameters and, at this dimensionality, we measure an acceptance fraction of only
$\sim$0.1. This is the practical reason hierarchical inference was deferred in Paper~I as too
costly, and it is precisely the barrier the differentiable forward model removes. Amortized SBI
remains orders of magnitude faster per evaluation once trained, but the hierarchical HMC requires
no training simulations and yields calibrated, interpretable population hyperparameters directly.

\subsection{Amortized hierarchical SBI: a unified approach}
\label{sec:amortized_hbi}

The two methods compared above need not be alternatives; they can be composed. As a proof of
concept we trained a single neural posterior estimator on \emph{individual} cluster profiles,
$q_\phi(M,c \mid d_j)$, amortizing the per-cluster likelihood. Because the training prior is
uniform, $q_\phi$ is proportional to the per-cluster likelihood and can be inserted directly into
the hierarchical model of Section~\ref{sec:hbm} in place of the analytic forward-model likelihood:
the per-cluster latents $(M_j,c_j)$ are drawn from the population $(\mu_M,\sigma_M,c_0,\beta,
\sigma_c)$, and each cluster contributes a likelihood term given by its amortized SBI posterior. We
summarize each per-cluster SBI posterior by a Gaussian and sample the population hyperparameters
with NUTS.

This ``amortized hierarchical SBI'' recovers the population \emph{mass} distribution well
($\mu_M=14.36$, $\sigma_M=0.13$ vs.\ true $14.38, 0.118$) in $\sim$10\,s of population sampling
after the one-time training, combining SBI's amortization with the hierarchical layer's coherent,
interpretable population inference.

The population \emph{concentration} spread, by contrast, is not recovered ($\sigma_c\approx0.47$
vs.\ true $0.18$), and we have verified that this is a fundamental information limit rather than a
methodological artifact. The per-cluster SBI posteriors are well calibrated (the standardized
residuals between the per-cluster posterior mean and the true value have unit width for both $M$
and $c$), but each single-cluster concentration is constrained only to $\sigma_c^{\rm meas}\approx
0.9$ at our fiducial $0.3$\,dex profile noise---roughly five times the true population spread. The
population scatter is therefore less than $4\%$ of the total observed concentration variance, below
the noise floor. Fixing the $M$--$c$ slope to the Child18 value or tightening the population
hyperprior does not change this ($\sigma_c$ remains $\approx0.47$), confirming that the limitation
is the per-cluster information content, not the prior or the slope. The same difficulty appears,
less severely, for hierarchical HMC (Table~\ref{tab:three_method}). This is a concrete statement of
what stacked weak lensing at this noise level can constrain: the population mass distribution is
recoverable, but the intrinsic concentration scatter is not, for any of the inference methods we
tested.

We regard amortized hierarchical SBI as the most promising direction emerging from this work: it is
fast, requires no differentiable forward model at inference time, and yields the same interpretable
population posterior as hierarchical HMC. \todo{Replace the per-cluster Gaussian summary with a
normalizing-flow per-cluster likelihood that preserves the non-Gaussian $M$--$c$ degeneracy, and
quantify the noise level at which the concentration spread becomes recoverable.}

\section{Discussion and conclusions}
\label{sec:conc}

We have shown that a hierarchical Bayesian model, enabled by a differentiable forward model and
gradient-based sampling, recovers the intrinsic population distribution of cluster mass and
concentration from stacked weak lensing --- closing the gap with SBI and demonstrating that the
``overconfidence'' of standard MCMC in Paper~I is a property of the single-$(M,c)$ assumption, not
of likelihood-based inference. The two approaches are complementary: SBI is amortized and
automatically robust to noise-induced biases, but requires representative training simulations and
its accuracy degrades under distribution shift; hierarchical HMC requires a differentiable model
and per-fit sampling, but needs no training set, provides direct and interpretable population
hyperparameters with built-in convergence diagnostics, and remains accurate under the scatter
regimes we tested.

One limitation of the present implementation is the assumption of a Gaussian population in
$\log_{10}M$. The true in-bin mass distribution, set by the steeply falling mass function
convolved with richness-selection scatter, is Gaussian-consistent in the baseline configuration
(Shapiro--Wilk $p=0.12$) but becomes measurably non-Gaussian (flat-topped, negative excess
kurtosis; $p<10^{-3}$) once the richness--mass scatter is large. In that regime the Gaussian
population model is mis-specified, which likely accounts for the slight discrepancy in the
recovered spread for the high $\lambda$--M scatter experiment. We deliberately do \emph{not} impose
the mass-function$\times$richness selection as a fixed prior on the population (as the
non-hierarchical informed-prior MCMC of Paper~I does to match SBI's training distribution), because
the goal of the hierarchical model is to \emph{infer} the population rather than assume it; a
natural refinement is a more flexible population family (e.g.\ skew-normal, or a
mass-function-informed shape with free amplitude and width).

\subsection{Toward hierarchical SBI}
\label{sec:hier_sbi}

The SBI approach used here and in Paper~I infers a fixed set of population summary statistics
(seven percentiles of the marginal $M$ and $c$ distributions plus their correlation) and then fits
a two-dimensional Gaussian to them. This is effective but has three structural limitations that
parallel the population-model issue above. First, the final Gaussian cannot represent a
non-Gaussian population; this is the likely cause of SBI's underestimate of the spread in the high
$\lambda$--M scatter experiment (Table~\ref{tab:three_method}), the same regime in which the true
in-bin mass distribution is significantly non-Gaussian. Second, compressing the population to a
handful of percentiles discards information and introduces estimator noise at finite $N_c$. Third,
inferring summaries and then fitting them is a two-stage procedure that does not propagate
uncertainty coherently.

Several more principled approaches exist. (i) One can train a neural posterior estimator to map the
\emph{set} of cluster profiles directly to the population hyperparameters $(\mu_M, \sigma_M, c_0,
\beta, \sigma_c)$, using a permutation-invariant set embedding (e.g.\ deep sets), yielding a single
coherent population posterior rather than a percentile fit. (ii) The fitted Gaussian can be
replaced by a flexible conditional density (a normalizing flow or mixture density network), which
removes the Gaussian assumption and would directly address the non-Gaussian regime. (iii) Most
compellingly, neural likelihood/ratio estimation can be used to \emph{amortize the per-cluster
likelihood}, which is then inserted into exactly the hierarchical model developed in this paper ---
``amortized hierarchical SBI.'' This last option unifies the two threads of this work: SBI provides
the fast, simulation-trained likelihood surrogate, while the hierarchical layer provides the
interpretable, correctly-propagated population inference. Dedicated simulation-based hierarchical
inference methods have been developed in the broader literature and would be a natural foundation.
We implement option~(iii) as a proof of concept in Section~\ref{sec:amortized_hbi}.
\todo{Option (iv): evaluate dedicated simulation-based hierarchical inference methods (e.g.\ the
hierarchical neural posterior estimation of Rozet \& Louppe 2021) as an alternative to the
amortized-likelihood approach used here.}

Several further extensions are natural: jointly inferring the richness--mass relation and its
scatter as hyperparameters; propagating photometric-redshift and miscentering systematics through
the hierarchy; and a head-to-head calibration comparison on survey-realistic mocks. We defer these
to future work. \todo{quantify HMC concentration bias mitigation; add multi-richness-bin joint fit;
add coverage/calibration test for the hierarchical posteriors; cite specific hierarchical-SBI
methods (e.g.\ Rozet \& Louppe 2021).}

\section*{Data and code availability}
All code, configurations, and the hierarchical-inference notebook are available in the
\code{CL-SBI} repository. \todo{add DOI/Zenodo link.}

\section*{Acknowledgments}
\todo{standard LSST-DESC acknowledgments.}

\bibliographystyle{plainnat}
\bibliography{hbi_paper}

\end{document}
